\chapter{Fundamentação Teórica}
\label{chap:fundamentacao}

A fundamentação teórica é onde você traz conceitos importantes, que servem como base para o entendimento do seu trabalho. Você traz autores consagrados, gráficos que você achou no Google Scholar às 2h da manhã, sempre citando a fonte \enquote{De acordo com os estudos de \cite{Ding2020AdversarialAttacks}}. 

Assim, vai construindo um texto onde seus parágrafos são retirados diretamente de diversos trabalhos, mostrando que você não tá inventando nada.

O segredo aqui é parecer que você mergulhou fundo na literatura, mesmo que tenha mergulhado mesmo só na aba de recomendados do Google Scholar.

Para isso, é ideal é criar vários subtópicos e seções. Você também pode incluir equações:

\input{equacoes/bhaskara}

Ou até mesmo código, onde você tem mais liberdade para trazer conhecimento elaborado por si mesmo.

\lstinputlisting[language=Python, 
caption=Exemplo de código
,label=lst:exemplocodigo]{chapters/trechos_codigo/python.m}
\hfill
\begin{minipage}[t]{.65\textwidth}
\fonteref{Elaborado pelo autor.}
\end{minipage}
%exemplo de inputs, ideal para organização e troca de posicionamento futuro é ter um elemento por arquivo
